% Generated draft: requires mathematical and source review.
\documentclass[11pt]{article}
\usepackage[margin=25mm]{geometry}
\usepackage{fontspec}
\setmainfont{Latin Modern Roman}
\usepackage{amsmath,amssymb,mathtools,unicode-math}
\setmathfont{Latin Modern Math}
\usepackage{enumitem,needspace,xurl,hyperref}
\setlength{\parindent}{0pt}
\setlength{\parskip}{5pt}
\setlength{\emergencystretch}{3em}
\newcommand{\N}{\mathbb{N}}
\newcommand{\Z}{\mathbb{Z}}
\newcommand{\Q}{\mathbb{Q}}
\newcommand{\R}{\mathbb{R}}
\newcommand{\C}{\mathbb{C}}
\newcommand{\F}{\mathbb{F}}
\newcommand{\A}{\mathbb{A}}
\newcommand{\PP}{\mathbb{P}}
\newcommand{\E}{\mathbb{E}}
\newcommand{\eps}{\varepsilon}
\newcommand{\dd}{\,\mathrm{d}}
\begin{document}
\section*{UnsolvedMath ID: 30000034}
\textbf{Source problem code:} catalogue entry
\subsection{Definitions and mathematical statement}
The source requests a classification of algebraic surfaces of general type with small invariants, explicitly including geometric genus zero and one. In the usual complex smooth projective setting, a surface $S$ is of general type when its Kodaira dimension is $\kappa(S)=2$. Its basic invariants are
\[
 p_g(S)=h^0(S,K_S),\qquad q(S)=h^1(S,\mathcal O_S),\qquad
 \chi(\mathcal O_S)=1-q(S)+p_g(S),\qquad K_S^2.
\]
Here $K_S$ is the canonical divisor and $K_S^2$ its self-intersection. A minimal smooth surface contains no smooth rational curve of self-intersection $-1$; using minimal models is necessary if $K_S^2$ is to enter a birational classification without arbitrary blowups.

The requested program is to determine the relevant surfaces and their families for small allowed values of these invariants, in particular for
\[
 p_g(S)=0\qquad\hbox{and}\qquad p_g(S)=1.
\]
The source does not give numerical bounds for ``small'', does not state whether these two cases exhaust its intended range, and does not fix the other invariants. It also does not specify classification up to isomorphism, birational equivalence, or deformation, nor whether to use smooth minimal surfaces or singular canonical models. The complex smooth projective definitions above explain the customary framework; the coefficient field, regularity and model convention require confirmation rather than being silently imposed on the broad request.

The supplied literature assessment checked 21 August 2026 says that a complete classification with small invariants is unavailable and that even the $p_g=0$ and $p_g=1$ programs remain incomplete. This dated, broad assessment does not say that every fixed numerical subcase is open and is not an independent current literature finding.
\subsection{Short English statement}
Classify general-type surfaces with small numerical invariants, including geometric genus zero and one. A precise classification target still needs bounds on the other invariants, a model convention and an equivalence relation. The supplied August 2026 assessment reports the overall program incomplete.
\subsection{Sources}
\begin{itemize}
\item Local UnsolvedMath registry, permanent ID 30000034, exact source code catalogue entry; \emph{Classification of General-Type Surfaces with Small Invariants}. Proposed year 2004, status open, exact supplied wording.
\item \emph{Mini-Workshop: Classification of Surfaces of General Type with Small Invariants (2004)}, Oberwolfach Reports, DOI 10.4171/owr/2004/09, pp. 439--478, \url{https://doi.org/10.4171/owr/2004/09}.
\item Supplied literature triage checked 2026-08-21: ``Oberwolfach report'', \url{https://ems.press/journals/owr/articles/719}; ``Surfaces with p\_g=0'', \url{https://arxiv.org/abs/1004.2583}. This dated OWR-derived curation requests independent verification of time-sensitive status and cited claims.
\end{itemize}
\end{document}
